\section{Traceability and scope limits}

\subsection{Client requirements to functional requirements}

\begin{center}
\small
\begin{tabularx}{\textwidth}{@{}L{0.08\textwidth}L{0.52\textwidth}Y@{}}
\toprule
\textbf{Client} & \textbf{Requirement} & \textbf{Functional requirements} \\
\midrule
CR1 & Yearly tax statement per donor & FR-D6, FR-P1 \\
CR2 & Validate volunteer skill and availability before assignment & FR-V3 \\
CR3 & Certificate after the required cumulative hours & FR-V4, FR-V5, FR-V6 \\
CR4 & Notify attendees when an event is updated or rescheduled & FR-E5, FR-C2 \\
CR5 & Record aid distribution per beneficiary & FR-B2, FR-B3, FR-P1 \\
CR6 & Full CRUD for event types, and for skills & FR-R1, FR-R2, FR-R3, FR-R4 \\
CR7 & Email, SMS, and social channels, swappable at runtime & FR-C1, FR-C2 \\
\bottomrule
\end{tabularx}
\end{center}

\subsection{Out of scope for this version}

The following items are deliberately excluded from Phase~1. They are listed so that the boundary of the deliverable is clear:

\begin{itemize}
  \item user accounts, login, and role-based access control;
  \item real payment processing and real message delivery;
  \item multi-branch operation and concurrent multi-user access;
  \item mobile or web clients;
  \item the Phase~2 design patterns and the admin pages.
\end{itemize}

\subsection{Acceptance of the phase}

The phase is accepted when the Must requirements in Section~3 can be demonstrated from the running application, the reference-data module performs all four operations on event types and skills, and the six Phase~1 patterns are present in the class diagram and in the code.
